\documentclass[11pt,a4paper]{article}
\usepackage[margin=2cm]{geometry}
\usepackage[T1]{fontenc}
\usepackage{mathptmx}
\usepackage{amsmath,amssymb}
\usepackage{xcolor}
\usepackage{enumitem}
\usepackage{array}
\usepackage[hidelinks]{hyperref}
\definecolor{navy}{HTML}{1B3755}
\definecolor{hl}{HTML}{FFF2CC}
\setlength{\parindent}{0pt}
\setlength{\parskip}{5pt}
\renewcommand{\arraystretch}{1.22}
\newcommand{\sect}[1]{\bigskip{\Large\color{navy}\bfseries #1}\par\smallskip{\color{navy}\hrule height 0.8pt}\smallskip}
\newcommand{\say}[1]{\par\smallskip\noindent\fbox{\parbox{\dimexpr\linewidth-2\fboxsep-2\fboxrule}{\textbf{In one sentence:} #1}}\par}

\begin{document}
\begin{center}
{\LARGE\bfseries\color{navy} The LG 3\,$\Omega$ Fault Table Explained}\\[4pt]
{\large Sheet \texttt{LG\_3ohm\_faults} of \texttt{Minimum\_Fault\_Current\_Calculation.xlsx}}\\[2pt]
IEEE 13-node feeder, DIgSILENT PowerFactory, DG out of service
\end{center}

\sect{1. What the table is}
Each row is one short-circuit calculation in PowerFactory: a \textbf{line-to-ground (LG) fault through a fault
resistance of 3\,$\Omega$} at one node, on one phase, with the DG disconnected. This is the \emph{weakest} fault of the
study, the one that gives the smallest fault current. It is used to check that each recloser's pickup is low enough
to see every fault of its zone. The calculation uses the \emph{complete method}, i.e.\ it starts from the load flow,
so the loads keep drawing current during the fault.

\sect{2. The columns}
\begin{center}
\begin{tabular}{|l|>{\raggedright\arraybackslash}p{12.2cm}|}
\hline
\textbf{Column} & \textbf{Meaning} \\ \hline
Faulted node & Node where the fault is applied (DL = the distributed load on line 632--671). \\ \hline
Faulted phase & Phase that is short-circuited to ground. Only the phases that exist at the node appear: 645 and 646 have
phases B and C, 684 has A and C, 652 only A, 611 only C. \\ \hline
Rf (ohm) & Fault resistance, 3\,$\Omega$ for every row. \\ \hline
Current at the fault (A) & Current flowing into the fault itself, i.e.\ to ground at the faulted node. \\ \hline
Current through R1 (A) & Largest phase current measured at R1 (feeder head) during the fault: fault current \emph{plus}
the load current that R1 still carries. \\ \hline
Current through R2 (A) & Largest phase current measured at R2 (on line 632--671) during the fault. \\ \hline
In R1 zone & ``yes'' if the fault is in the zone R1 must protect (every node except 634). \\ \hline
In R2 zone & ``yes'' if the fault is downstream of R2 (671, 692, 675, 680, 684, 652, 611). \\ \hline
Yellow cells & The lowest recloser current of each zone: this is \textbf{If,min} used in the pickup check. \\ \hline
\end{tabular}
\end{center}

\sect{3. How to read the numbers}
\textbf{(a) Why is the R1 current larger than the fault current?}\\
Because R1 carries everything that flows into the feeder: the fault current \emph{and} the load current of the phase.
Example, row 680 phase B: the fault draws 701.5\,A, but R1 measures 1074.4\,A. With phasors:
$I_{R1} \approx I_{load,R1} + I_{fault}$ = 411\,A\,$\angle{-173.5^\circ}$ + 701.5\,A\,$\angle{-176.0^\circ}$
$\approx$ 1112\,A (PowerFactory: 1074\,A, a little less because the loads draw less while the voltage sags).

\textbf{(b) Why does R2 show about 465\,A for faults at 645, 646 or DL?}\\
These faults are \emph{upstream} of R2 (between the feeder head and R2). With the DG out, no fault current flows through
R2; R2 only carries the load beyond it (about 465--470\,A). So the ``In R2 zone'' column is ``no'' for them and these rows
are not used for R2.

\textbf{(c) Why are rows 671 and 692 identical?}\\
671 and 692 are joined by a switch of zero length, so electrically they are the same point.

\textbf{(d) Why is node 634 excluded?}\\
634 is on the 0.48\,kV side of transformer XFM-1. A 3\,$\Omega$ fault there draws only about 91\,A at 4.16\,kV --- about the
transformer's own load current. That node is protected by its own fuse on the transformer, not by R1 or R2, so it is
not part of either recloser zone for this check.

\textbf{(e) Why is the lowest value at 680 and not at the farthest node 652?}\\
The farthest node (652) gives the smallest current \emph{at the fault} (617.2\,A). But the recloser current also contains
the load current of the faulted phase. Phase B carries the least load, so the LG fault on phase B at 680 gives the
smallest current \emph{through the recloser}: 1074.4\,A at R1 and 878.5\,A at R2.

\sect{4. What the table proves}
\begin{center}
\begin{tabular}{|l|c|c|c|c|}
\hline
\textbf{Recloser} & \textbf{If,min (yellow)} & \textbf{Where} & \textbf{Pickup used} & \textbf{Check} \\ \hline
R1 & 1074.4\,A & 680, phase B & 720\,A & 720 $<$ 1074 \checkmark \\ \hline
R2 & 878.5\,A & 680, phase B & 600\,A (curve start) & 600 $<$ 879 \checkmark \\ \hline
\end{tabular}
\end{center}
Both pickups are below the weakest fault current of their zones, so both reclosers detect every fault they must
protect against (without DG).

\say{The table lists the weakest fault (LG through 3\,$\Omega$) at every node and phase; for each recloser the smallest
current it measures for a fault in its own zone is 1074\,A (R1) and 879\,A (R2), both above the pickups of 720\,A and
600\,A, so the pickups are correctly set.}

\sect{5. ``Isn't the fault a sink that takes everything?''}
Only for a bolted fault. With 3\,$\Omega$ the fault does not short the feeder out:

\textbf{(1) The voltage stays up.} At the fault point $V \approx I_f \times R_f = 701.5\,\text{A} \times 3\,\Omega
\approx 2.1$\,kV, against 2.58\,kV before the fault. The feeder still runs at about 80\,\% voltage, so every load keeps
drawing most of its normal current.

\textbf{(2) Kirchhoff's current law at R1.} With the DG out, R1 is the only entry from the grid, so everything that
flows anywhere on the feeder passes through it:
\[
I_{R1} = I_{fault} + \sum I_{loads}
\]
For the fault at 680 (phase B): 701.5\,A into the fault plus 411\,A of phase-B load. The two are almost in phase
($-176^\circ$ and $-173.5^\circ$), so they nearly add: $\approx$ 1112\,A. PowerFactory gives 1074\,A, a little less
because the loads draw somewhat less at 80\,\% voltage.

\textbf{(3) When the fault does take almost everything.} For a bolted fault (0\,$\Omega$) the voltage near the fault
collapses and the loads draw very little, so R1 carries about the fault current alone. Example: bolted 3-phase fault
at 633, 4102\,A at the fault and 4162\,A at R1 --- only about 60\,A apart.

\textbf{Why it matters.} The relay measures the total current of its phase, load plus fault. The pickup must
therefore be compared with what the relay actually sees, 1074\,A for R1 and 879\,A for R2.

\sect{6. Likely questions}
\textbf{Why 3\,$\Omega$?} It is the fault resistance prescribed by the method for the minimum fault (weakest
realistic fault).\\
\textbf{Why the largest phase current at the recloser?} A phase-overcurrent relay operates on the phase with the
largest current; that is the current it compares with its pickup.\\
\textbf{What changes with DG?} With the DG in service the grid share falls; at 100\,\% DG the weakest fault drops to
577\,A at R1 and 368\,A at R2, below the pickups (shown in the DG penetration study).
\end{document}
